\newcommand{\PraxisTitle}{How much can averaging reduce bounded error?}
% BEGIN PRAXIS RESEARCH STYLE v2
% Praxis mathematical research house style v2; target-journal instructions take precedence.
\documentclass[11pt,a4paper]{amsart}
\usepackage[T1]{fontenc}
\usepackage{lmodern}
\usepackage{amsmath,amssymb,amsthm,booktabs,caption,needspace,placeins}
\usepackage[margin=30mm]{geometry}
\usepackage[hidelinks,pdfauthor={},pdftitle={\PraxisTitle}]{hyperref}
\urlstyle{same}
\usepackage{pgfplots}
\pgfplotsset{compat=1.17}
\usepgfplotslibrary{fillbetween,groupplots}
\usetikzlibrary{arrows.meta,positioning,calc}
\definecolor{main}{HTML}{0072B2}
\definecolor{accent}{HTML}{D55E00}
\definecolor{muted}{HTML}{7A7F85}
\definecolor{light}{HTML}{B8BDC2}
\definecolor{oigreen}{HTML}{009E73}
\definecolor{oiorange}{HTML}{E69F00}
\definecolor{oisky}{HTML}{56B4E9}
\definecolor{oiyellow}{HTML}{F0E442}
\definecolor{teal}{HTML}{17766E}
\definecolor{clay}{HTML}{9E6440}
\definecolor{slate}{HTML}{5D696E}
\definecolor{sand}{HTML}{AAA69B}
\captionsetup{font=small,labelfont=bf,labelsep=period,justification=raggedright,singlelinecheck=true}
\captionsetup[table]{position=above,skip=5pt}
\captionsetup[figure]{position=below,skip=7pt}
\raggedbottom
\renewcommand{\topfraction}{.9}\renewcommand{\textfraction}{.08}\renewcommand{\floatpagefraction}{.8}
\setlength{\emergencystretch}{2em}
\newtheorem{theorem}{Theorem}[section]
\newtheorem{lemma}[theorem]{Lemma}
\newtheorem{proposition}[theorem]{Proposition}
\newtheorem{corollary}[theorem]{Corollary}
\theoremstyle{definition}
\newtheorem{definition}[theorem]{Definition}
\theoremstyle{remark}
\newtheorem{remark}[theorem]{Remark}
\renewcommand{\proofname}{Proof}

% Localize labels only when ctex is loaded; English documents retain AMS defaults.
\usepackage{etoolbox}
\makeatletter
\newcommand{\PraxisResearchChinese}{%
  \renewcommand{\abstractname}{摘要}%
  \renewcommand{\keywordsname}{关键词}%
  \renewcommand{\datename}{日期}%
  \renewcommand{\refname}{参考文献}%
  \renewcommand{\proofname}{证明}%
  \patchcmd{\abstract}{\scshape\abstractname.}{\normalfont\bfseries\abstractname\quad}{}%
    {\PackageWarning{praxis}{Chinese abstract label patch did not apply}}%
  \def\@setdate{\datename：\enspace\@date}%
  \def\@setkeywords{\keywordsname：\enspace\@keywords}%
}
\AtBeginDocument{\@ifpackageloaded{ctex}{\PraxisResearchChinese}{}}
\makeatother
% END PRAXIS RESEARCH STYLE

% Working-paper example. Replace the question, evidence, metadata and references.
\title[Bounded error in a weighted mean]{\PraxisTitle}
\author{}\date{}
\keywords{weighted mean, worst-case error}
\begin{document}
\begin{abstract}
Averaging measurements with a common deterministic error bound leaves
the worst-case error unchanged. For any nonnegative weights that sum to
one, the bound is attained when every error equals the common upper bound.
\end{abstract}
\maketitle
\section{Measurement model}
Let $y_i=\theta+e_i$ for $1\leq i\leq n$, where $n\geq1$ and
$|e_i|\leq\delta$. The errors have no assumed probability law.
We ask how far a weighted estimate can lie from $\theta$.
\section{A sharp bound}
\begin{proposition}
If $w_i\geq0$, $\sum_i w_i=1$ and $\delta\geq0$, then
\[
  \sup_{|e_i|\leq\delta}\left|\sum_i w_i y_i-\theta\right|=\delta .
\]
\end{proposition}
\begin{proof}
The error is $\sum_i w_i e_i$. Its absolute value is at most
$\sum_i w_i\delta=\delta$, and setting every $e_i=\delta$ attains equality.
\end{proof}
The supremum is a dual-norm calculation \cite{bv}. Adding observations
alone cannot reduce this bound. A claim about variance would require a
probability model for the errors.
\begin{thebibliography}{1}
\bibitem{bv} S. Boyd and L. Vandenberghe, \emph{Convex Optimization},
Cambridge University Press, 2004, Appendix A.
\end{thebibliography}
\end{document}
